\documentclass[12pt,oneside,a4paper,chapter=TITLE]{abntex2}

\usepackage[T1]{fontenc}
\usepackage[utf8]{inputenc}
\usepackage[brazil]{babel}
\usepackage{indentfirst}
\usepackage{graphicx}
\usepackage{lmodern}
\usepackage{hyperref}
\usepackage{float}

\renewcommand{\thesection}{\arabic{section}}
\renewcommand{\theHsection}{\arabic{section}}
\setcounter{section}{0}

\titulo{TempleOS: Análise de uma Arquitetura de Sistema Operacional de Usuário Único}

\autor{
Natan Gabriel de Souza Miranda \\
João Victor Soares de Lima \\
Kamilly Stephany Saraiva Santos \\
Samuel Alexandre Correia da Silva \\
Crytian Pinheiro Aguiar \\
Roberto Matheus Nunes dos Santos \\
Paulo Israel Pinho Filho
}

\instituicao{Centro Universitário Maurício de Nassau - UNINASSAU}
\local{Maceió - AL}
\data{2026}
\tipotrabalho{Relatório Técnico}

\begin{document}

\hypersetup{
    colorlinks=true,
    linkcolor=blue,
    citecolor=black,
    urlcolor=blue
}

\begin{figure}[h]
\centering

\includegraphics[width=4cm]{Logo_unina.png}
\end{figure}

\imprimircapa

\begin{folhaderosto}

\begin{center}

\includegraphics[width=4cm]{Logo_unina.png}
\vspace{2cm}

{\ABNTEXchapterfont\Large\imprimirautor}

\vfill

{\ABNTEXchapterfont\bfseries\LARGE\imprimirtitulo}

\vfill
\end{center}

\hspace*{8cm}
\begin{minipage}{8cm}
Trabalho apresentado à disciplina Sistemas Operacionais do curso de Ciência da Computação do Centro Universitário Maurício de Nassau --- UNINASSAU, como requisito parcial para obtenção da nota da AV2.

\vspace{0.5cm}
Professor(a): Anderlan Henrique
\end{minipage}

\vfill

\begin{center}
\imprimirlocal \\
\imprimirdata
\end{center}

\end{folhaderosto}

\tableofcontents
\newpage

\textual

\section{Resumo}

O TempleOS é um sistema operacional experimental desenvolvido por Terry A. Davis entre 2003 e 2018. Ele se destaca por sua simplicidade extrema, arquitetura de usuário único e ausência de rede. Este trabalho analisa suas características técnicas, linguagem própria (HolyC), limitações e execução prática em ambiente virtualizado.

\textbf{Palavras-chave:} TempleOS; Sistemas Operacionais; HolyC; Arquitetura Computacional.

\section{Introdução}

O TempleOS é um sistema operacional experimental projetado com foco em simplicidade extrema e acesso direto ao hardware.

Diferente de sistemas modernos, ele elimina camadas de segurança e abstração, permitindo interação direta com o sistema.

Este trabalho apresenta uma análise técnica do sistema.

\section{Contexto Histórico}

O TempleOS foi desenvolvido por Terry A. Davis ao longo de aproximadamente 15 anos, sendo finalizado por volta de 2013.

Davis foi responsável por todas as partes do sistema, incluindo kernel, interface gráfica e linguagem de programação.

\begin{figure}[H]
\centering
\includegraphics[width=5cm]{Terry_Davis.jpg}
\caption{Terry A. Davis, criador do TempleOS}
\end{figure}

\section{Arquitetura do Sistema}

O TempleOS possui arquitetura de 64 bits, single-user e single-core.

Principais características:
\begin{itemize}
    \item Kernel monolítico simples;
    \item Sem permissões de usuário;
    \item Execução direta de código;
    \item Interface gráfica própria;
    \item Sistema altamente integrado.
\end{itemize}

\section{Linguagem HolyC}

O sistema utiliza a linguagem HolyC, derivada do C.

Características:
\begin{itemize}
    \item Compilação em tempo real;
    \item Integração com o sistema;
    \item Programação interativa;
    \item Acesso direto à memória.
\end{itemize}

\begin{figure}[H]
\centering
\includegraphics[width=0.9\textwidth]{HOLYC.JPG}
\caption{Exemplo de código em HolyC no ambiente do TempleOS}
\end{figure}

\section{Segurança e Limitações}

Limitações do sistema:
\begin{itemize}
    \item Sem suporte a rede;
    \item Sem isolamento de processos;
    \item Sem autenticação de usuários;
    \item Sem proteção de memória.
\end{itemize}

\section{Interface do Sistema}

Interface simples com resolução fixa de 640x480 pixels e 16 cores.

\section{Interface e Elementos Visuais do Sistema}

\begin{figure}[H]
\centering
\includegraphics[width=0.9\textwidth]{Menu.jpg}
\caption{Menu do sistema}
\end{figure}

\begin{figure}[H]
\centering
\includegraphics[width=0.9\textwidth]{Tela_inicial.jpg}
\caption{Tela inicial}
\end{figure}

\begin{figure}[H]
\centering
\includegraphics[width=0.9\textwidth]{Passagem_biblica.jpeg}
\caption{Passagem bíblica}
\end{figure}

\begin{figure}[H]
\centering
\includegraphics[width=0.9\textwidth]{Keymap.jpg}
\caption{Keymap do sistema}
\end{figure}

\section{Ambiente de Testes e Execução em Máquina Virtual}

O TempleOS foi executado em ambiente virtualizado utilizando Oracle VM VirtualBox.

Observações:
\begin{itemize}
    \item Inicialização rápida;
    \item Interface limitada;
    \item Forte integração entre sistema e linguagem;
    \item Ausência de recursos modernos.
\end{itemize}

\section{Unicidade do TempleOS: A Visão de Um Criador}

O TempleOS é fundamentalmente único por ter sido desenvolvido por uma única pessoa --- Terry A. Davis --- que manteve total controle sobre o design e implementação de todos os componentes do sistema.

Essa característica distingue o TempleOS de outros sistemas operacionais modernos:

\begin{itemize}
    \item \textbf{Coesão de Design:} Ao contrário de sistemas desenvolvidos por grandes equipes (como Linux ou Windows), o TempleOS reflete uma visão singular e coerente de arquitetura de sistema operacional;
    \item \textbf{Simplicidade Intencional:} Davis eliminou deliberadamente complexidades desnecessárias, criando um sistema altamente integrado onde cada componente foi projetado para trabalhar em harmonia;
    \item \textbf{Linguagem Própria:} A criação da linguagem HolyC especificamente para o TempleOS permite uma integração profunda entre linguagem e sistema que seria difícil de alcançar em projetos colaborativos;
    \item \textbf{Documentação e Intencionalidade:} Cada decisão arquitetural foi tomada com propósito definido, resultando em um sistema altamente coeso onde nenhuma parte foi adicionada por pressão externa ou requisitos de compatibilidade;
    \item \textbf{Liberdade Criativa:} A ausência de comitês, padrões industriais ou restrições comerciais permitiu que Davis explorasse conceitos inovadores e não convencionais.
\end{itemize}

Esta unicidade torna o TempleOS um estudo de caso valioso para entender como sistemas operacionais podem ser arquitetados quando princípios de simplicidade e coesão são prioritários.

\section{Comparação com Sistemas Operacionais Modernos}

O TempleOS diferencia-se significativamente dos sistemas operacionais modernos em vários aspectos fundamentais:

\begin{itemize}
    \item \textbf{Arquitetura Multi-usuário vs. Single-usuário:} Sistemas modernos como Linux, Windows e macOS suportam múltiplos usuários simultâneos com diferentes níveis de permissão, enquanto TempleOS é estritamente single-user;
    \item \textbf{Segurança e Isolamento:} Sistemas modernos implementam mecanismos robustos de proteção de memória, sandboxing de processos e controle de acesso. O TempleOS elimina essas proteções em favor da simplicidade;
    \item \textbf{Conectividade de Rede:} Os sistemas operacionais modernos incluem suporte nativo para redes e conectividade à internet, funcionalidades completamente ausentes no TempleOS;
    \item \textbf{Compatibilidade de Hardware:} Enquanto sistemas modernos funcionam em diversas arquiteturas (x86, ARM, PowerPC, etc.), o TempleOS foi desenvolvido especificamente para arquitetura x86 64-bits;
    \item \textbf{Modelo de Compilação:} Sistemas modernos utilizam compiladores que geram código otimizado e separado. O TempleOS adota compilação Just-In-Time (JIT) em tempo real como parte integral do sistema;
    \item \textbf{Gerenciamento de Recursos:} Sistemas modernos oferecem escalonamento sofisticado de processos, gerenciamento de memória virtual e swapping. O TempleOS executa tudo em um único núcleo sem complexidades de escalonamento;
    \item \textbf{Tamanho e Complexidade:} O kernel do TempleOS contém aproximadamente 65 mil linhas de código em HolyC, enquanto kernels modernos (Linux, Windows) possuem milhões de linhas;
    \item \textbf{Filosofia de Design:} Sistemas modernos priorizam versatilidade, compatibilidade e recursos avançados. O TempleOS prioriza minimalismo, clareza e integração total entre componentes.
\end{itemize}

Essa comparação demonstra que o TempleOS representa uma abordagem radicalmente diferente ao design de sistemas operacionais, servindo como contraste educacional para compreender as decisões arquitetônicas dos sistemas modernos.

\section{Impacto e Relevância}

O TempleOS possui relevância acadêmica por seu design minimalista e abordagem não convencional.

É estudado em arquitetura de sistemas operacionais e design de software.

\section{Conclusão}

O TempleOS demonstra uma abordagem única de sistema operacional, priorizando simplicidade extrema e integração total entre software e hardware.

A análise em máquina virtual reforça suas limitações e características experimentais.
\begin{thebibliography}{99}

\bibitem{templeos}
TempleOS Official. Disponível em: https://templeos.org/. Acesso em: 2026.

\bibitem{wiki_templeos}
WIKIPEDIA. TempleOS. Disponível em: \url{https://en.wikipedia.org/wiki/TempleOS}. Acesso em: 2026.

\bibitem{tanenbaum}
TANENBAUM, Andrew S.; BOS, Herbert. \textit{Modern Operating Systems}. Pearson, 2015.

\bibitem{ostep}
ARPACI-DUSSEAU, Remzi; ARPACI-DUSSEAU, Andrea. \textit{Operating Systems: Three Easy Pieces}. 2018.

\end{thebibliography}

\end{document}